\documentclass[11pt,a4paper]{article}

\usepackage{amsmath,amssymb,amsthm}
\usepackage[margin=1in]{geometry}
\usepackage{booktabs}
\usepackage{hyperref}
\usepackage{enumitem}

\hypersetup{
    colorlinks=true,
    linkcolor=black,
    urlcolor=blue,
    citecolor=black
}

\newcommand{\vp}{\varphi}
\newcommand{\vpbar}{\bar{\varphi}}
\newcommand{\Mpl}{M_{\mathrm{Pl}}}
\newcommand{\me}{m_e}
\newcommand{\mpr}{m_p}
\newcommand{\GeV}{\,\mathrm{GeV}}

\title{Interface Theory: Core Equations and Verification}
\author{S.\ Kittilsen}
\date{}

\begin{document}
\maketitle

\begin{abstract}
All modular forms are evaluated at a single algebraically distinguished nome
$q = 1/\vp$, where $\vp = (1+\sqrt{5})/2$ is the golden ratio.
This document presents the resulting formulas for Standard Model coupling constants
and other physical quantities, the logical chain connecting them, uniqueness tests,
and falsifiable predictions. All computations are reproducible via standard Python
scripts with no dependencies beyond the \texttt{math} module.
\end{abstract}

\tableofcontents

%----------------------------------------------------------------------
\section{Setup: Modular Forms at the Golden Nome}
\label{sec:setup}

\subsection{Definitions}

The Dedekind eta function and Jacobi theta functions are defined by their standard
$q$-series:
\begin{align}
  \eta(q) &= q^{1/24} \prod_{n=1}^{\infty} (1 - q^n), \\[6pt]
  \theta_2(q) &= 2 q^{1/4} \sum_{n=0}^{\infty} q^{n(n+1)}, \\[6pt]
  \theta_3(q) &= 1 + 2\sum_{n=1}^{\infty} q^{n^2}, \\[6pt]
  \theta_4(q) &= 1 + 2\sum_{n=1}^{\infty} (-1)^n\, q^{n^2}.
\end{align}

\subsection{The Golden Nome}

All expressions in this document are evaluated at:
\begin{equation}
  \boxed{q = \frac{1}{\vp} = \frac{\sqrt{5}-1}{2} = 0.6180339887\ldots}
\end{equation}
This is the unique positive solution of $q + q^2 = 1$.

\subsection{Numerical Values}

\begin{center}
\begin{tabular}{lll}
\toprule
Function & Value at $q = 1/\vp$ & Digits used \\
\midrule
$\eta(1/\vp)$ & $0.118404\ldots$ & \\
$\theta_2(1/\vp)$ & $2.55509\ldots$ & $\approx \theta_3$ \\
$\theta_3(1/\vp)$ & $2.55509\ldots$ & \\
$\theta_4(1/\vp)$ & $0.03031\ldots$ & \\
$\vp$ & $1.6180339887\ldots$ & \\
$\vpbar \equiv 1/\vp$ & $0.6180339887\ldots$ & $= \vp - 1$ \\
\bottomrule
\end{tabular}
\end{center}

%----------------------------------------------------------------------
\section{The Derivation Chain}
\label{sec:chain}

The logical structure proceeds through seven steps.  Each step is listed with
its key equation or theorem and a brief justification.

\medskip
\noindent\textbf{Step 0.} \quad $q + q^2 = 1 \implies q = 1/\vp$.\\
The unique positive root of the self-referential equation.  The iteration
$f(x) = 1/(1+x)$ converges to $1/\vp$ from any positive starting value.

\medskip
\noindent\textbf{Step 1.} \quad The $E_8$ root lattice contains $\mathbb{Z}[\vp]^4$.\\
$E_8$ is the unique exceptional simple Lie algebra whose root lattice embeds the
golden ring $\mathbb{Z}[\vp]$.  The minimal polynomial of $\vp$ is
$x^2 - x - 1 = 0$, i.e.\ $q + q^2 = 1$.

\medskip
\noindent\textbf{Step 2.} \quad The unique quartic potential in $\mathbb{Z}[\vp]$ is:
\begin{equation}
  \boxed{V(\Phi) = \lambda\bigl(\Phi^2 - \Phi - 1\bigr)^2.}
\end{equation}
Two degenerate vacua at $\Phi = \vp$ and $\Phi = -1/\vp$.  This is the only
self-consistent quartic in $E_8$'s golden field
\cite{derive_V_from_E8}.

\medskip
\noindent\textbf{Step 3.} \quad The kink solution
$\Phi_{\mathrm{kink}}(x) = \frac{\sqrt{5}}{2}\tanh(\kappa x) + \frac{1}{2}$
interpolates between the two vacua.  The fluctuation operator yields a
P\"oschl--Teller potential of depth $n = 2$:
\begin{equation}
  V_{\mathrm{fluct}}(x) = -\frac{6\kappa^2}{\cosh^2(\kappa x)},
\end{equation}
which has exactly two bound states and is \emph{reflectionless}
($|T(k)|^2 = 1$ for all $k$).  This is a standard result in
soliton physics~\cite{Rajaraman1982}.

\medskip
\noindent\textbf{Step 4.} \quad The Lam\'e equation on the kink lattice gives nome
$q = 1/\vp$ via the elliptic modulus $k(1/\vp) = 0.9999999901\ldots \to 1$
(P\"oschl--Teller limit).  Spectral invariants of the Lam\'e operator
evaluated at this nome produce modular forms at $q = 1/\vp$.

\medskip
\noindent\textbf{Step 5.} \quad SM coupling constants emerge as ratios of these
modular forms (Section~\ref{sec:couplings}).

\medskip
\noindent\textbf{Step 6.} \quad Nome doubling: Jacobi nome $q_J = 1/\vp$ vs.\
modular nome $q_M = 1/\vp^2$.  Three couplings exhaust the ring of
$\Gamma(2)$-modular forms.  Verified: the creation identity
$\eta(q)^2 = \eta(q^2)\cdot\theta_4(q)$ holds to $3.7\times 10^{-16}$.

%----------------------------------------------------------------------
\section{Coupling Constants}
\label{sec:couplings}

\subsection{Fine Structure Constant $\alpha$}

\subsubsection{Tree Level}

\begin{equation}
  \frac{1}{\alpha_0} = \frac{\theta_3(q)}{\theta_4(q)}\cdot\vp
  = \frac{2.55509\ldots}{0.03031\ldots}\times 1.61803\ldots
  = 136.39\ldots
\end{equation}

\subsubsection{VP-Corrected (Self-Consistent Fixed Point)}

The full formula includes vacuum polarization running with a halved coefficient,
as required by the Jackiw--Rebbi theorem~\cite{JackiwRebbi1976}
(chiral zero mode on a domain wall contributes one Weyl fermion,
not one Dirac fermion, to the VP loop):
\begin{equation}
  \boxed{
  \frac{1}{\alpha} = \frac{\theta_3}{\theta_4}\cdot\vp
  \;+\; \frac{1}{3\pi}\ln\!\left(\frac{\Lambda_{\mathrm{ref}}}{\me}\right)
  }
\end{equation}
where
\begin{align}
  \Lambda_{\mathrm{ref}} &= \frac{\mpr}{\vp^3}\,f(x), &
  x &= \frac{\eta}{3\vp^3}, \\[6pt]
  f(x) &= 1 - x + \tfrac{2}{5}\,x^2.
\end{align}
The coefficient $c_2 = 2/5 = n/(2n+1)$ for P\"oschl--Teller depth $n = 2$ is
derived from the Graham--Weigel kink pressure integral~\cite{GrahamWeigel2024}.

The closed-form VP function is the confluent hypergeometric (Kummer) function:
\begin{equation}
  f(x) = \tfrac{3}{2}\;{}_1F_1(1;\,\tfrac{3}{2};\,x) - 2x - \tfrac{1}{2}.
\end{equation}

\paragraph{Result.}
\begin{center}
\begin{tabular}{ll}
\toprule
Predicted & $1/\alpha = 137.0359991$ \\
Measured (CODATA 2018) & $1/\alpha = 137.035999084(21)$ \\
\midrule
Precision & 10.2 significant figures \\
Residual & 0.062\,ppb, $\;0.4\sigma$ \\
\bottomrule
\end{tabular}
\end{center}

\subsubsection{Supporting Derivations}

\begin{itemize}[nosep]
\item VP coefficient $1/(3\pi)$ instead of $2/(3\pi)$:
      Jackiw--Rebbi~\cite{JackiwRebbi1976} $+$ Kaplan mechanism~\cite{Kaplan1992}
      $+$ APS index theorem $+$ Callan--Harvey anomaly inflow~\cite{CallanHarvey1985}.
\item $c_2 = 2/5$: Wallis integral ratio from kink one-loop pressure,
      Graham \& Weigel (PLB 852, 2024).
\item Self-consistent fixed point: the wall counts its own bound states via
      $\zeta_{\mathrm{bs}}(0) = n = 2$ (spectral zeta function) and adjusts
      its coupling accordingly.
\end{itemize}

\subsection{Strong Coupling $\alpha_s$}

\begin{equation}
  \boxed{\alpha_s = \eta(1/\vp) = 0.11840}
\end{equation}

\begin{center}
\begin{tabular}{ll}
\toprule
Predicted & $0.11840$ \\
Measured (PDG 2024) & $0.1179 \pm 0.0009$ \\
Deviation & $0.6\sigma$ \\
\midrule
\textbf{Live test} & \textbf{CODATA 2026--27} \\
\bottomrule
\end{tabular}
\end{center}

This is a committed prediction with zero free parameters: the strong coupling
\emph{is} the Dedekind eta function at the golden nome.

\subsection{Weinberg Angle $\sin^2\!\theta_W$}

\begin{equation}
  \boxed{\sin^2\!\theta_W
  = \frac{\eta^2}{2\,\theta_4} - \frac{\eta^4}{4}
  = 0.23126}
\end{equation}

\begin{center}
\begin{tabular}{ll}
\toprule
Predicted & $0.23126$ \\
Measured & $0.23122 \pm 0.00004$ \\
Deviation & $0.3\sigma$ \\
\bottomrule
\end{tabular}
\end{center}

\subsection{Core Identity}

The three coupling-related quantities satisfy:
\begin{equation}
  \boxed{
  \alpha^{3/2}\,\mu\,\vp^2
  \left[1 + \frac{\alpha\ln\vp}{\pi}\right] = 3
  }
\end{equation}
where $\mu = \mpr/\me = 1836.15267\ldots$ is the proton-to-electron mass ratio.
Tree level ($99.89\%$); with $1$-loop correction ($99.999\%$, a $122\times$ improvement).
The correction $\alpha\ln(\vp)/\pi$ has the form of a one-loop gauge
correction in the kink background with vacuum ratio~$\vp^2$.

%----------------------------------------------------------------------
\section{Additional Quantities}
\label{sec:additional}

\subsection{Cosmology and Hierarchy}

\begin{center}
\begin{tabular}{@{}llllr@{}}
\toprule
Quantity & Formula & Predicted & Measured & Match \\
\midrule
$\Lambda/\Mpl^4$ & $\theta_4^{80}\sqrt{5}/\vp^2$ & $2.88\times10^{-122}$
  & $2.89\times10^{-122}$ & $\sim$exact \\[4pt]
$v$ (Higgs VEV) & $\Mpl\,\vpbar^{80}/(1-\vp\theta_4)\cdot(1+\eta\theta_4\cdot7/6)$
  & $246.218\GeV$ & $246.22\GeV$ & 99.999\% \\[4pt]
$\mu\;(\mpr/\me)$ & $6^5/\vp^3 + 9/(7\vp^2)$
  & $1836.156$ & $1836.15267$ & 99.9998\% \\[4pt]
$m_t$ (top) & $\me\cdot\mu^2/10$
  & $172.57\GeV$ & $172.69\GeV$ & 99.93\% \\[4pt]
$\eta_B$ (baryon asymmetry) & $\theta_4^6/\!\sqrt{\vp}$
  & $6.098\times10^{-10}$ & $6.12\times10^{-10}$ & 99.6\% \\[4pt]
$\Omega_m/\Omega_\Lambda$ & $\eta(1/\vp^2)$
  & $0.4625$ & $0.4599$ & 99.4\% \\[4pt]
$\gamma_{\mathrm{Immirzi}}$ & $1/(3\vp^2)$
  & $0.12738$ & $0.12732$ & 99.95\% \\
\bottomrule
\end{tabular}
\end{center}

\noindent
The exponent $80$ appearing in $\Lambda$ and $v$ is the Coxeter sum through the
exceptional chain:
\begin{equation}
  h(E_8) + h(E_7) + h(E_6) + h(D_5) + h(A_4) + h(A_3) + h(A_2)
  = 30 + 18 + 12 + 8 + 5 + 4 + 3 = 80.
\end{equation}

\subsection{Fermion Masses (Zero Free Parameters)}

All nine charged-fermion masses are expressed as $m = \mpr\times(\text{formula})$,
using only $\{\vp,\,\mu,\,3,\,4/3,\,10,\,2/3\}$:

\begin{center}
\begin{tabular}{@{}llrrr@{}}
\toprule
Fermion & $m/\mpr$ & Predicted & Measured & Error \\
\midrule
$e$   & $1/\mu$ & 0.000545 & 0.000545 & 0.00\% \\
$u$   & $\vp^3/\mu$ & 0.002307 & 0.002302 & 0.21\% \\
$d$   & $9/\mu$ & 0.004902 & 0.004977 & 1.52\% \\
$\mu$ (muon) & $1/9$ & 0.11111 & 0.11261 & 1.33\% \\
$s$   & $1/10$ & 0.10000 & 0.09954 & 0.46\% \\
$c$   & $4/3$ & 1.33333 & 1.35355 & 1.49\% \\
$\tau$ & Koide($e,\mu$), $K\!=\!2/3$ & 1.89387 & 1.89376 & 0.01\% \\
$b$   & $4\vp^{5/2}/3$ & 4.44025 & 4.45500 & 0.33\% \\
$t$   & $\mu/10$ & 183.615 & 184.051 & 0.24\% \\
\bottomrule
\end{tabular}
\end{center}

\noindent
Average error: $0.62\%$.\quad Maximum error: $1.52\%$.\quad Free parameters: $0$.

The assignment rule uses the $S_3$ representation action on $\mathbb{Z}/4\mathbb{Z}$
via the Chinese Remainder Theorem isomorphism
$\mathbb{Z}/12\mathbb{Z} \cong \mathbb{Z}/3\mathbb{Z}\times\mathbb{Z}/4\mathbb{Z}$.
The flavor symmetry $S_3 = \mathrm{SL}(2,\mathbb{Z})/\Gamma(2)$ is proven
mathematics~\cite{Feruglio2017}.

\subsection{Mixing Matrices}

\subsubsection{PMNS (Neutrino Mixing)}

\begin{center}
\begin{tabular}{@{}llllr@{}}
\toprule
Angle & Formula & Predicted & Measured & $\sigma$ \\
\midrule
$\sin^2\!\theta_{12}$ & $1/3 - \theta_4\sqrt{3/4}$
  & $0.3071$ & $0.3092 \pm 0.0087$ & $0.24\sigma$ \\
$\sin^2\!\theta_{23}$ & $1/2 + 40C$\;\footnotemark
  & $0.5718$ & $0.572$ & $<0.1\sigma$ \\
$\sin^2\!\theta_{13}$ & $3\,\theta_4/(\theta_3\,\vp)$
  & $0.02200$ & $0.02203 \pm 0.00056$ & $0.06\sigma$ \\
\bottomrule
\end{tabular}
\end{center}
\footnotetext{$C = \eta\,\theta_4/2 = 0.001794\ldots$ is a universal loop correction
factor that also appears in the alpha and Higgs VEV formulas.}

\noindent
$\sin^2\!\theta_{12} = 0.3071$ is a \textbf{live experimental test} (JUNO, ongoing).

\subsubsection{CKM (Quark Mixing)}

\begin{center}
\begin{tabular}{@{}lllr@{}}
\toprule
Element & Formula & Predicted / Measured & Match \\
\midrule
$|V_{us}|$ & $(\vp/7)(1-\theta_4)$ & $0.2241$ / $0.2243$ & 99.5\% \\
$|V_{cb}|$ & $(\vp/7)\sqrt{\theta_4}$ & $0.0402$ / $0.0405$ & 99.4\% \\
$|V_{ub}|$ & $(\vp/7)\cdot 3\,\theta_4^{3/2}(1+\vp\theta_4)$
  & $0.00384$ / $0.00382$ & 99.5\% \\
$\delta_{CP}$ & $\arctan\!\bigl(\vp^2(1-\theta_4)\bigr)$
  & $68.5^\circ$ / $68.4^\circ$ & 99.9997\% \\
\bottomrule
\end{tabular}
\end{center}

%----------------------------------------------------------------------
\section{Uniqueness Tests}
\label{sec:uniqueness}

\subsection{Algebra Uniqueness}

Among the five exceptional Lie algebras $\{G_2, F_4, E_6, E_7, E_8\}$,
\textbf{only $E_8$} produces formulas matching all three coupling constants
($\alpha$, $\alpha_s$, $\sin^2\!\theta_W$).  The others achieve $0/3$.
The discriminating property: only $\mathbb{Z}[\vp]$ (the golden ring, embedded
in $E_8$) has discriminant $+5$ (positive $\Rightarrow$ real quadratic field
$\Rightarrow$ two real vacua $\Rightarrow$ domain wall exists).  Other
exceptional algebras have negative discriminant (complex roots, no domain wall).

\emph{Verification:} \texttt{lie\_algebra\_uniqueness.py}

\subsection{Nome Uniqueness}

A scan of $6{,}061$ rational and algebraic nomes $q \in (0,1)$ finds that
$q = 1/\vp$ is the \textbf{only} nome for which all three formulas
($\alpha_s$, $\sin^2\!\theta_W$, $1/\alpha$ tree level) simultaneously match
measured values within $1\%$.

\emph{Verification:} \texttt{nome\_uniqueness\_scan.py}

\subsection{Formula Isolation}

$719$ neighboring formulas (obtained by systematic perturbation of exponents,
prefactors, and functional forms) were tested against the core identity.
\textbf{Zero} match at $1\%$.  The formula is structurally isolated in
formula space.

\emph{Verification:} \texttt{formula\_isolation\_test.py}

\subsection{Statistical Significance}

Monte Carlo testing (see \texttt{expected\_match\_count.py}):
matching $\sim\!30$ individual quantities to $<1\%$ is not by itself surprising
given the functional freedom available.  The genuinely significant result is that
\textbf{three core formulas} (the three couplings) simultaneously match at $<1\%$
using a single nome---this occurs with probability $0.2\%$ in the null model.
The $10.2$-significant-figure match for $\alpha$ goes far beyond what any
random search produces.

%----------------------------------------------------------------------
\section{Falsification}
\label{sec:falsification}

\subsection{Four Live Experimental Tests}

\begin{center}
\begin{tabular}{@{}lllll@{}}
\toprule
Test & Prediction & Current status & Instrument & Timeline \\
\midrule
$\alpha_s$ & $0.11840$ & $0.1179 \pm 0.0009$ ($0.6\sigma$)
  & Lattice QCD / CODATA & 2026--27 \\
$\sin^2\!\theta_{12}$ & $0.3071$ & $0.3092 \pm 0.0087$ ($0.24\sigma$)
  & JUNO & ongoing \\
$r$ (tensor/scalar) & $0.0033$ & $< 0.036$
  & CMB-S4 / LiteBIRD & $\sim$2028 \\
$R = d\!\ln\mu/d\!\ln\alpha$ & $-3/2$ & untested
  & ELT / ANDES & $\sim$2035 \\
\bottomrule
\end{tabular}
\end{center}

The \textbf{decisive test} is $R = -3/2$: if confirmed, the likelihood ratio
exceeds $10^{10}$.  The \textbf{sharpest near-term blade} is $\alpha_s = 0.11840$
(CODATA 2026--27).

\subsection{Dead Claims (19 Removed)}

The framework has killed its own wrong ideas.  Removed items include:

\begin{enumerate}[nosep]
\item $\pi = \theta_3(1/\vp)^2\cdot\ln\vp$ -- tautological (holds for any large $q$).
\item Coupling triangle $\alpha_s^2 = 2\sin^2\!\theta_W\cdot\theta_4$ -- tautological
      (substituting definitions gives $\eta^2 = \eta^2$).
\item CKM unitarity -- tautological (built into parameterization).
\item Creation identity -- Jacobi triple product rearranged; true for all $q$.
\item Muon $g-2$ -- wrong sign on higher-order coefficients.
\item $H_0$ from Lucas/Fibonacci -- decorative (cherry-picked indices).
\item $\theta_2 \approx \theta_3$ -- generic for large $q$.
\item Jacobi abstrusa -- classical identity, all $q$.
\item Lam\'e gap ratio $= 3$ -- P\"oschl--Teller limit, not $\vp$-specific.
\item Harris current sheet $n = 1$ -- topological, not tunable.
\item Solar $p$-mode golden ratios -- negative result (no $\vp$ concentration).
\item VP second derivation via 2D hexagonal lattice -- honest negative.
\item Bare $V(\Phi)$ inflation -- fails ($\eta = -3.2$ at hilltop).
\item Fermion decay rate $c = \ln\vp$ -- false (actual $c = 2$).
\item All 9 masses via Kaplan mechanism -- original fails (6 free params).
\item $\theta_4 \equiv 1\pmod{p}$ fingerprint -- null test failed ($p = 0.686$).
\item $196560/196884$ GUT unification -- numerology.
\item $L(2,\chi_5)/\vp^3 = 1/6$ -- impossible (Lindemann--Weierstrass).
\item $\zeta_K(2)/\pi^2 = \eta(1/\vp)$ -- $2.5\sigma$ gap, not exact.
\end{enumerate}

%----------------------------------------------------------------------
\section{Verification}
\label{sec:verification}

All scripts are standard Python~3, requiring only the built-in \texttt{math} module
(no NumPy, no SciPy).  Modular form values are computed from truncated $q$-series
(convergence is rapid at $q = 1/\vp \approx 0.618$; 50 terms suffice for
$>15$-digit precision).

\bigskip
\noindent
Repository: \url{https://github.com/kittilsenstian-debug/the-hand}

\begin{center}
\begin{tabular}{@{}ll@{}}
\toprule
Script & Verifies \\
\midrule
\texttt{verify\_golden\_node.py} & All golden nome numerical claims \\
\texttt{modular\_couplings\_v2.py} & $\alpha_s$, $\sin^2\!\theta_W$, tree-level $\alpha$ \\
\texttt{alpha\_self\_consistent.py} & $\alpha$ at 10.2 significant figures \\
\texttt{alpha\_cascade\_closed\_form.py} & VP closed form via ${}_1F_1$ \\
\texttt{derive\_V\_from\_E8.py} & $V(\Phi)$ uniqueness from $E_8$ \\
\texttt{lie\_algebra\_uniqueness.py} & $E_8$ unique ($3/3$ vs $0/3$) \\
\texttt{nome\_uniqueness\_scan.py} & $q = 1/\vp$ unique ($6{,}061$ tested) \\
\texttt{formula\_isolation\_test.py} & Core identity isolated ($0/719$) \\
\texttt{expected\_match\_count.py} & Monte Carlo null model \\
\texttt{one\_resonance\_fermion\_derivation.py}
  & 9 fermion masses, 0 free parameters \\
\texttt{kink\_1loop.py} & 1-loop correction $\alpha\ln(\vp)/\pi$ \\
\texttt{exponent\_80\_completion.py} & Coxeter sum $= 80$ \\
\texttt{three\_generations\_derived.py} & 3 generations from $\Gamma(2) \to S_3$ \\
\texttt{coupling\_triangle\_sigma.py} & Coupling triangle with measured values \\
\bottomrule
\end{tabular}
\end{center}

%----------------------------------------------------------------------
\begin{thebibliography}{99}

\bibitem{JackiwRebbi1976}
R.~Jackiw and C.~Rebbi,
``Solitons with fermion number 1/2,''
\emph{Phys.\ Rev.\ D} \textbf{13}, 3398 (1976).

\bibitem{Kaplan1992}
D.~B.~Kaplan,
``A method for simulating chiral fermions on the lattice,''
\emph{Phys.\ Lett.\ B} \textbf{288}, 342 (1992).

\bibitem{CallanHarvey1985}
C.~G.~Callan and J.~A.~Harvey,
``Anomalies and fermion zero modes on strings and domain walls,''
\emph{Nucl.\ Phys.\ B} \textbf{250}, 427 (1985).

\bibitem{GrahamWeigel2024}
N.~Graham and H.~Weigel,
``Kink one-loop pressure,''
\emph{Phys.\ Lett.\ B} \textbf{852}, 138607 (2024).

\bibitem{Feruglio2017}
F.~Feruglio,
``Are neutrino masses modular forms?''
in \emph{From My Vast Repertoire: Guido Altarelli's Legacy},
A.~Ferrara, ed., World Scientific (2017).

\bibitem{RubakovShaposhnikov1983}
V.~A.~Rubakov and M.~E.~Shaposhnikov,
``Do we live inside a domain wall?''
\emph{Phys.\ Lett.\ B} \textbf{125}, 136 (1983).

\bibitem{Rajaraman1982}
R.~Rajaraman,
\emph{Solitons and Instantons}, North-Holland (1982).

\bibitem{RandallSundrum1999}
L.~Randall and R.~Sundrum,
``A large mass hierarchy from a small extra dimension,''
\emph{Phys.\ Rev.\ Lett.} \textbf{83}, 3370 (1999).

\bibitem{Borcherds1992}
R.~E.~Borcherds,
``Monstrous moonshine and monstrous Lie superalgebras,''
\emph{Invent.\ Math.} \textbf{109}, 405 (1992).

\bibitem{derive_V_from_E8}
Verification script: \texttt{derive\_V\_from\_E8.py},
available at \url{https://github.com/kittilsenstian-debug/the-hand}.

\end{thebibliography}

\end{document}
